\section{1. Introduction}

Light field imaging has emerged as a powerful modality in modern computer vision, capturing both spatial and angular information of light rays, as demonstrated in foundational light field imaging and depth estimation studies . This rich multidimensional data facilitates accurate depth estimation, which is essential for downstream applications such as autonomous navigation, robotic manipulation, and immersive 3D reconstruction, as highlighted in recent applied vision surveys . By analyzing the epipolar plane image (EPI)s, conventional algorithms extract geometric disparities based on the linear trajectories of light rays across different viewpoints. Consequently, light field depth estimation achieves notable performance in controlled environments, establishing a solid foundation for complex spatial perception tasks. These geometric cues provide a distinct advantage over traditional monocular or stereo vision techniques, especially in textureless regions.

However, obtaining the precise geometry of non-Lambertian objects remains a substantial challenge, since standard photo-consistency cues cannot be easily applied to complex reflectances . Most current depth estimation methods primarily support the Lambertian assumption, making them unreliable for glossy, metallic, or translucent surfaces . In reality, light interacting with complex materials undergoes diverse physical processes, including Fresnel specular reflection, Rayleigh or Mie scattering, and cosine diffuse reflection. Such situations break the photo-consistency assumptions, leading to erroneous depth estimation results in most existing algorithms . Therefore, the inherent physical complexity of non-Lambertian reflectance severely degrades the reliability of purely geometric light field models.

The severity of this performance degradation is evident in empirical evaluations. For instance, a standard epipolar plane image (EPI) baseline model achieves a mean absolute error of 0.081 in mixed urban scenes, but the error deteriorates significantly to 0.411 in purely non-Lambertian domains. This failure occurs because the network predictions degenerate into mere texture replication rather than structural depth recovery. Furthermore, identifying material properties from limited angular samples is highly difficult. A conventional black-box three-class convolutional neural network yields a mere 24.6\\% accuracy when classifying reflectance types from $9 \times 9$ angular patches. Additionally, the extreme scarcity of non-Lambertian training data, with only four available scenes, further restricts the generalization capability of data-driven approaches.

Addressing these limitations is critically necessary for deploying light field systems in unconstrained real-world environments. Relying solely on implicit feature extraction without physical constraints inevitably leads to ambiguous representations in mixed-material scenes. A unified framework must explicitly model the light-matter interaction to decouple different reflection components. By quantifying surface roughness and angular deflection vectors, the system can dynamically adapt its depth estimation strategy for diffuse, specular, and scattering regions. Consequently, developing a physically grounded model that unifies Lambertian and non-Lambertian depth estimation is not only theoretically significant but also practically urgent for advancing robust 3D vision systems.

To address these challenges, we review existing depth estimation frameworks, which can be broadly categorized into traditional geometric approaches, early deep learning models, and recent advanced networks.

Traditional light field depth estimation methods primarily rely on epipolar plane image (EPI)s to extract geometric disparities, as proposed in early geometric frameworks . However, these approaches operate on the strict assumption of global Lambertian reflectance. In non-Lambertian regions, specular reflections and scattering alter the angular distribution of light, causing non-linear trajectories in epipolar plane image (EPI)s . Consequently, these geometric models lack explicit modeling of light-matter interactions, resulting in poor robustness and frequent failures in mixed-material scenes.

To mitigate the limitations of geometric cues, early deep learning models utilize convolutional neural networks to implicitly extract features from stacked viewpoints, exemplified by architectures like EPINet . Despite their improved performance on Lambertian surfaces, these data-driven approaches suffer from critical methodological flaws. They typically treat material parsing as a black-box classification task and rely on entirely implicit feature extraction. Consequently, these networks struggle to spontaneously learn physically meaningful representations, which restricts their generalization under complex illumination and significantly reduces model interpretability.

More recent approaches attempt to address these issues by introducing complex architectures, such as defocus branches or pre-trained backbone networks, as seen in recent advanced models . While these methods aim to incorporate additional cues, they encounter significant technical bottlenecks, such as ineffective depth feature learning and severe overfitting on relatively small light field datasets. Moreover, these advanced networks typically employ a unified depth estimation strategy across the entire image. This uniform processing inherently restricts their ability to adaptively switch estimation strategies based on pixel-level reflection mechanisms, highlighting an urgent need for a physics-driven, unified approach.

In this paper, we propose a unified dual-mask physical model for non-Lambertian light field depth estimation. Unlike previous data-driven approaches that rely on implicit feature extraction, our framework explicitly integrates physical optics principles into the network architecture. Specifically, we leverage angular frequency domain analysis to parse material properties and employ a dual-mask mechanism to modulate light field features. This physical formulation effectively decouples complex light-matter interactions, enabling robust depth recovery across Lambertian, specular, and scattering regions within a single unified model.

Our main contributions are summarized as follows:

\begin{itemize}
  \item \textbf{Contribution 1}: We propose an angular frequency material parsing mechanism inspired by magnetic resonance imaging k-space analysis. By applying a two-dimensional discrete Fourier transform to the $9 \times 9$ angular sampling matrix, we establish a physical mapping between angular spatial frequencies and material reflectance properties, enabling accurate parsing of Lambertian, specular, and scattering regions.
  \item \textbf{Contribution 2}: We design a physics-driven dual-mask feature modulation module that explicitly decouples complex light-matter interactions. This mechanism dynamically adapts the feature extraction process based on the parsed material properties, ensuring robust representation learning across diverse reflection components.
  \item \textbf{Contribution 3}: We introduce an adaptive multi-strategy depth estimation framework that seamlessly integrates distinct depth recovery pathways for different material types. Extensive experiments on challenging non-Lambertian light field datasets demonstrate that our unified model significantly outperforms state-of-the-art methods in both accuracy and robustness.
\end{itemize}